\documentclass[12pt]{article}
\usepackage[utf8]{inputenc}
\usepackage[T1]{fontenc}
\usepackage[margin=1in]{geometry}
\usepackage{setspace}
\usepackage{xcolor}
\usepackage[colorlinks=true, urlcolor=blue!60!black, linkcolor=blue!60!black]{hyperref}
\setlength{\parindent}{0em}
\setlength{\parskip}{0.6em}
\renewcommand{\baselinestretch}{1.0}
\begin{document}
\begin{center}
{\LARGE \textit{Hayburn's Case}, 2 U.S. (2 Dall.) 409 (1792)\par}
\vspace{0.5em}
{\large Case Brief --- Two-Page Edition\par}
\end{center}

\section{Memory Jogger}
Five Justices on circuit refused to decide veterans' pension claims subject to revision by the Secretary of War and Congress: courts may perform only judicial duties, and their decisions must be final (\href{https://supreme.justia.com/cases/federal/us/2/409/}{Justia}).

\section{Detailed Case Facts}
The Invalid Pensions Act of 1792 had circuit courts determine disabled veterans' eligibility and certify findings to the Secretary of War, who could withhold a name on ``cause to suspect imposition or mistake''; Congress could also revise the awards.

\section{Procedural History}
The Pennsylvania Circuit refused to hear William Hayburn's petition, and three circuits wrote to President Washington. Attorney General Randolph sought mandamus; the Supreme Court split 3--3 on his \textit{ex officio} authority, held the renewed motion under advisement, and Congress mooted the case in 1793.

\section{Judicial Votes}
No Supreme Court merits judgment. Jay, Cushing, Wilson, Blair, and Iredell stated their views in circuit letters.

\section{Holding}
Through the letters, five Justices established that Congress may not assign Article III courts non-judicial duties or subject their judgments to executive or legislative revision.

\section{Analysis of Opinions}
\textbf{The circuit letters:} The three branches are distinct and independent; courts may be given only duties ``properly judicial'' and performed ``in a judicial manner.'' Revision by the Secretary of War or Congress would make courts subordinate advisers. The circuits differed on whether judges could serve individually as commissioners: New York did so, Pennsylvania refused, and North Carolina doubted it. \emph{Later:} \textit{Plaut} (1995) barred reopening final judgments; \textit{Muskrat} (1911) applied the advisory-opinion bar; \textit{Mistretta} (1989) used the commissioner theory.

\section{Examples: Future Applications}
\textbf{Same-Side:} a statute letting the Homeland Security Secretary vacate district-court awards is invalid. \textbf{Opposite-Side:} judges may serve voluntarily on a non-adjudicatory commission. \textbf{Fence-Sitter:} veterans' judgments payable only if Congress appropriates funds for each award.

\section{Critique}
The judges condemned advisory opinions in unsolicited letters to the President, and the commissioner device elevates form over independence. The only actual ruling was a 3--3 tie, later elevated from a reporter's footnote into canon.

\section{Key Quotations}
\begin{quote}
``That by the Constitution of the United States, the government thereof is divided into three distinct and independent branches, and that it is the duty of each to abstain from, and to oppose, encroachments on either. That neither the Legislative nor the Executive branches can constitutionally assign to the Judicial any duties but such as are properly judicial, and to be performed in a judicial manner.''
\end{quote}
\noindent --- New York Circuit Court (Jay, C.J., Cushing, J., Duane, D.J.), 2 U.S. (2 Dall.) at 410 n.*.

\end{document}
